\section{Introduction}
\label{sec:intro}

Diffusion models generate data by following a \emph{denoising trajectory}: starting from random noise, the model applies a sequence of learned noise-reduction steps until a sample from the target distribution emerges~\citep{ho2020denoising,song2019generative,song2021scorebased}. Generating one sample takes $T$ such denoising steps, each costing one forward pass through the network; we call each pass a \emph{function evaluation} (NFE), so standard sampling costs $T$ NFE.
Recent work shows that spending more than $T$ NFE per sample, a regime known as \emph{inference-time scaling}, can improve sample quality without retraining~\citep{ma2024inference,ramesh2025nts}.
This raises a resource-allocation question: how should the extra NFE be distributed across the $T$ denoising steps?
The extra budget can enter in two ways: (i)~\emph{lengthening the trajectory} by increasing $T$, which empirically saturates~\citep{ma2024inference}, and (ii)~\emph{searching within each timestep} by evaluating multiple noise candidates and keeping the one that scores highest on a quality metric.
Once the trajectory is long enough that (i) gives diminishing returns, additional NFE is best spent on (ii); we focus on this within-step search regime, which is the lever exploited by recent inference-time methods.

Within-step search is implemented as a local operation on the noise variable that drives each step.
Concretely, the model draws several candidate noise variables, scores each via a \emph{verifier} (a scalar quality function, e.g., a CLIP-style image-text alignment score), and keeps the highest-scoring one, obtaining a better sample at the cost of additional NFE.
Existing methods instantiate this local noise search in various forms (greedy, zero-order, and particle-based search; \citealt{ramesh2025nts,tang2024dno,kim2025rbf}).

These methods all enrich the local operator within a single step, but distribute the search budget uniformly across denoising steps or fix it to a hand-chosen window (e.g., only the initial noise, or a predetermined middle segment).
Yet replacing a noise draw with a better candidate empirically improves quality far more at some timesteps than at others; we call this per-step responsiveness the \emph{sensitivity} of step~$t$.
This leaves a complementary question: given a fixed NFE budget, \emph{where along the trajectory} should search concentrate?

We cast \emph{noise trajectory search} as a budget-allocation problem: choose per-step counts $K_t \in \mathbb{Z}_{\ge 1}$ that maximize the expected verifier score subject to $\sum_t K_t = B$. The problem decomposes into two levels coupled only through $K_t$: a \emph{local operator} searches noise candidates within one timestep, and a \emph{global scheduler} splits $B$ across the $T$ timesteps. The local level is well studied; the global level, which allocates differently across timesteps with unequal sensitivity, is the gap we address. This decoupling lets any existing local operator (greedy, zero-order, particle-based) compose with any global schedule; standard sampling, best-of-$N$, $\epsilon$-greedy, and particle/tree-based methods become special cases with hard-coded or uniform global policies.

Building on this formulation, we propose \textbf{\textsc{GAINS}} (\textbf{G}lobal \textbf{A}daptive \textbf{I}nference-time \textbf{N}oise \textbf{S}cheduling), a global scheduling algorithm that combines \emph{offline profiling} with \emph{online control}.
Offline, \textsc{GAINS} profiles verifier sensitivity across timesteps to obtain a coarse compute allocation specific to the model and dataset.
Online, \textsc{GAINS} adjusts the search at each step using the verifier score improvement and candidate variance observed so far, stopping early on unproductive steps and reallocating saved budget to later ones.
The controller operates causally (no access to future timesteps) and respects a strict NFE budget.

\paragraph{Contributions.}
\begin{enumerate}
\item We introduce a \emph{two-level view} of noise trajectory search that separates local noise refinement from global timestep scheduling, covering both stochastic and deterministic samplers and recovering existing inference-time methods as special-case scheduling policies (\cref{sec:framework}).
\item We develop \textsc{GAINS}, a \emph{global scheduling algorithm} that integrates offline timestep profiling with online feedback control, enforcing an exact NFE budget (\cref{sec:algorithm}).
\item We give a \emph{theoretical model} for the noise trajectory search problem with two optimality results: a closed-form water-filling characterization of the optimal offline allocation (\cref{thm:offline}), and a matching worst-case regret lower bound of order $T$ (\cref{lem:lower}) for a sequential-allocation game into which the diffusion problem embeds. The lower bound shows that no online-only policy avoids regret of order $T$ in the worst case, motivating the hybrid design.
\end{enumerate}

Empirically, on text-to-image (Stable Diffusion), class-conditional (EDM~\citep{karras2022elucidating}), and flow-based generation, \textsc{GAINS} matches or exceeds uniform-allocation benchmarks at fixed NFE budget, recovering the same Brightness and Compressibility verifier scores with 20--50\% fewer NFE, and adds $+0.005$ to $+0.04$ at equal budget; the gains hold across local search operators, prompt sets, and model families (\cref{sec:experiments}).

\subsection{Related Literature}
\label{sec:related-lit}

\paragraph{Inference-time noise search.}
A growing body of work improves diffusion sample quality by refining noise variables at generation time, paralleling LLM inference-time scaling~\citep{snell2024scaling,brown2024monkeys}.
Early approaches select among multiple full trajectories~\citep{ma2024inference} or optimize noise via gradient- or zero-order updates per timestep~\citep{tang2024dno,qi2024noises,zhou2025golden}.
More recent methods introduce richer local operators along three threads: greedy or two-phase search at each step~\citep{ramesh2025nts,kim2025rbf,zhang2025tscend,lee2025abcd,ttsnap2025}, particle- and tree-based search with non-uniform particle budgets across stages~\citep{singhal2025fk,kim2025psi,zhang2025classical}, and amortized or learned per-step compute~\citep{eyring2025hypernetworks,zhang2025adadiff,ye2025tpdm,yeh2025demons,sundaresha2025vt}.
The third thread modifies per-step effort but primarily adapts the base sampler rather than scheduling additional search compute.
These methods improve sample quality by designing richer \emph{local} search operators applied within each timestep. The complementary question of \emph{how much} budget each timestep should receive is typically resolved by uniform allocation or a hand-tuned schedule (e.g., particle counts fixed in advance in~\citet{singhal2025fk,kim2025psi}); we build on these local operators and address the global allocation question directly.
\textsc{GAINS} estimates a per-step sensitivity profile from a small calibration set and adapts allocation online based on observed verifier-score gains, distributing a fixed NFE budget non-uniformly across the trajectory.

\paragraph{Connections to bandits and resource allocation.}
Each component of the problem has a classical counterpart.
The local operator (evaluate $K$ noise candidates, keep the best) is a fixed-budget best-arm identification problem: split a sample budget across arms with unknown means and return the empirical best~\citep{chen2000simulation,kaufmann2016complexity,lattimore2020bandit}.
The global scheduler is a sequential budget-allocation problem across the $T$ steps~\citep{ibaraki1988resource,boyd2004convex}.
Our contribution to this perspective is a result specific to the diffusion setting (\cref{thm:loc-scale}) that decomposes the per-step expected gain into a sensitivity factor and a sample-size factor; this decomposition reduces the across-timestep allocation problem to a closed-form water-filling solution and gives a matching regret lower bound of order $T$.

\paragraph{Theoretical foundations of diffusion models.}
Several recent theoretical results provide useful context.
\citet{gao2024reward} develop a continuous-time RL framework for reward-directed diffusion and show that the optimal exploration policy has time-varying variance, a structure consistent with the location-scale model we adopt.
\citet{tang2024contractive} characterize a per-timestep error sensitivity for score-estimation propagation, conceptually related to our per-step sensitivity but capturing estimation rather than verifier-score sensitivity; \citet{chen2023sampling} establish polynomial convergence of score-based samplers under minimal data assumptions, foundational for the SDE discretization we use in \cref{sec:theory}.
Our work complements these analyses by focusing on the budget allocation problem: we show that per-step gains decompose into a sensitivity factor and a sample-size effect, and use this structure to characterize optimal scheduling.
